\section{Model Architectures}\label{sec:model-architectures}
In this project, a \gls{cnn} and an ensemble of \glspl{fb-mlp} were trained per patient to predict seizures.
This section discusses their exact specifications.

\subsection{Convolutional Neural Network}\label{subsec:cnn-architecture}
In this project, the architecture of the \gls{cnn} was adapted from \textcite{muller_coherent_2022}, which used the nv1x16 topology originally proposed in \textcite{eberlein_convolutional_2018}.
It analyzes a segment and uses the raw \gls{eeg} signals per \gls{eeg} channel as inputs.
Until the final stages, it treats \gls{eeg} channels separately, computing features per channel.
Only at the end are the per-channel features combined to output a final score.

In the following, \gls{eeg} channels refer to electrode-derived input signals, whereas \gls{cnn} channels refer to convolutional feature maps (tensor depth).

\Cref{fig:cnn-architecture} shows the exact architecture.
The inputs are passed through a combination of convolutional and pooling layers.
There are two distinct types of convolutional blocks: conv1 and conv2.
Conv1 uses a convolutional layer, followed by a pooling layer with a relatively large kernel size, reducing the height (\#samples) rapidly.
All convolution layers use L1 and L2 regularization with a penalty of $10^{-9}$.
Four conv1 blocks are followed by three conv2 blocks.
These further reduce the height by using valid padding for two convolutional layers and pooling again.
Ultimately, the sample dimension (height) is reduced to 1 with 32 features per \gls{eeg} channel.
Only then are the per-channel features combined (flattened).
These pass through two final dense layers to create the output logit, which is then bounded between 0 and 1 with the sigmoid function, to create the final output.

\begin{figure}[!htbp]
    \centering

    \begin{subfigure}[t]{\textwidth}
        \centering
        \includegraphics{conv_blocks}
        \caption{Special CNN Blocks}
        \label{fig:special-cnn-blocks}
    \end{subfigure}

    \vspace{1ex}

    \begin{subfigure}[t]{\textwidth}
        \centering
        \includegraphics[width=\linewidth]{CNN_architecture}
        \caption{Full CNN Architecture}
        \label{fig:full-cnn-architecture}
    \end{subfigure}

    \caption[CNN Architecture]{CNN Architecture.\\
    A segment's EEG data is processed from left to right. The volumes represent the input tensor and the outputs of layers. In between the volumes, the network's layer's properties are shown.\\
    \textit{k}: kernel size, \textit{f}: \#filters, \textit{p}: pool size.}
    \label{fig:cnn-architecture}
\end{figure}

Since the height is reduced further and further in deeper layers, the extracted features receive information from more samples, i.e., the receptive field of deeper neurons grows on the time scale.
Since the per-channel features are only combined after the final conv block, information from the \gls{eeg} channels is evaluated separately until the end.

To accommodate two differences in the data structure compared to \textcite{eberlein_convolutional_2018}, two adjustments were made to the architecture.
Here, 2 \gls{eeg} channels were used instead of 16.
This meant there were 8 times fewer nodes in the first flattened tensor.
For this reason, the number of nodes in the first dense layer was reduced by a factor of 8: from 64 to 8.
The flat tensors thus had a height of 64 → 8 → 1, instead of 512 → 64 → 1.

Secondly, because the sampling rate of $\sim$\SI{207}{\Hz} here differed slightly from the sampling rate of \SI{200}{\Hz}, the number of input samples from a segment differed too (3105 vs.\ 3000).
Thus, to arrive at a height of 1 at the final conv layer, it was necessary to reduce this dimension slightly more.
For this reason, the pool padding of the second conv2 layer was changed from same to valid.
This minor change was enough to fully accommodate the sampling frequency difference.

\subsection{Ensemble of Feature-Based Multilayer Perceptrons}\label{subsec:ensemble-architecture}
The second model was an ensemble of \glspl{fb-mlp} that was also adapted from \textcite{muller_coherent_2022}.
Instead of raw signal, it used hand-crafted features.
The input features consisted of the band power spectrum and correlation coefficient from \cref{subsec:feature-extraction}, which had previously demonstrated forecasting potential in a cycle-based framework \parencite{yang_seizure_2024}.
This allowed it to take features of individual channels and their correlation into consideration.

An individual \gls{fb-mlp} used the z-score normalized features as its input layer.
It had three dense, hidden layers comprising 16, 8, and 4 nodes, which used \gls{relu} as an activation function.
Batch normalization followed the input and hidden layers.
The output layer used the sigmoid activation function.

The ensemble consisted of 50 \glspl{fb-mlp} that were initialized with different random weights and shared their input layer.
The outputs of the individual models were then averaged to obtain the final output.
\Textcite{muller_coherent_2022} used an ensemble size of 100, rather than 50.
However, the original work investigated a broader set of input features, allowing for greater variability between models. 
Since the variability cannot be as high here, a smaller ensemble size was chosen to reduce resources needed for training.

\subsection{Model Training} \label{subsec:model-training}
For model training, only preictal and interictal segments were included. 
Inter-pre and intervention segments were excluded because they occur immediately prior to seizure onset and may therefore share electro-physiological characteristics with preictal segments. 
If included, these segments—while still annotated with the label 0 (non-preictal)—could induce the model to learn misleading patterns.

Similarly, postictal and inter-post segments were omitted, as they may fall within the preictal window of seizures that did not meet the validity criteria and could consequently contain preictal-like features. 
Moreover, the postictal period is characterized by substantial alterations in brain state, leading to marked changes in signal dynamics that may confound model learning and degrade performance \parencite{muller_coherent_2022}.

As mentioned in \cref{subsec:dataset-partitioning}, there was a large class imbalance with scarcity of preictal segments.
To avoid wasting resources training on more interictal segments than necessary, the number of interictal segments was subsampled to a maximum of 20 times the number of preictal ones.
The class weights were adjusted to the final ratio of preictal vs.\ interictal segments, so that the models were not biased towards focusing mostly on learning to predict interictal segments.
All models received the same subsampled segments for training, so that results were more comparable.

As in \textcite{eberlein_convolutional_2018}, the \glspl{cnn} were trained for 50 epochs with ADAM optimization, a learning rate of \num{0.001}, and binary cross-entropy as a loss function.

To train the \glspl{fb-mlp}, before subsampling, a z-score normalizer was fit to the features of the entire train set.
This same scaling was then always applied when this model was later used, including for the test set.
The \glspl{fb-mlp} were then trained for 200 epochs using stochastic gradient descent for optimization, a learning rate of \num{0.0001}, and binary cross-entropy as a loss function.
\Textcite{muller_coherent_2022} used 500 epochs for training.
However, here the performance could not be seen to improve after 200 epochs, which is likely also due to the smaller input dimension, requiring less fitting.
